

\section{Conclusions and outlook}\label{sec:Conclusions}

The present chapter has provided an overview of various types of
methods used in comparative corpus linguistics and translation and
interpreting studies. Some of the methods are well known, -- e.g.,
mosaic plots (as in Figure \ref{fig:mosaicplots1-2}) -- while others
may be less familiar to parts of the community of contrastive
linguists and translation
scholars.

We discussed two major types of ordination: CA (with and without
stacking), and MCA (with and without supplementary variables). We made
an effort to convey an intuitive understanding of the mathematics
underlying these techniques. Users have to make specific choices and
decisions, e.g., CA vs. MCA, the type of map (symmetric, asymmetric,
row principal, etc.), and such decisions need to be
informed.
%
We also highlighted the connection between (M)CA and the familiar
Chi-squared statistic (e.g., the connections between this statistic
and Chi-squared distance, inertia, and the $\phi$-coefficient).

It seems to us that graph-based methods have some potential that has
not fully been exploited. Context-Conditional Correlation Graphs
(CCCGs, cf. Sect. Fig. \ref{fig:CCCG} in
Sect. \ref{sec:LanguageSpecific}) allow us to represent a large amount
of information about pairwise associations between variables in a
single plot. Sankey diagrams (Figures \ref{fig:SankeyDeFr},
\ref{fig:SankeyDeFrRef} and \ref{fig:SankeyDeFrCluster}), which are
commonly used in other fields, do not seem to be widespread in
translation studies, though they offer a very good way of summarizing
directional data. Translation graphs visualize choices made in
translation, under specific conditions. While the type of graph that
we provided (Figure \ref{fig:NetworkDeFr} in
Sect. \ref{sec:ComparativeDirected}) does not contain any type of
statistical information, such information could be added, potentially
using Bayesian statistical reasoning as outlined in
Sect. \ref{sec:GraphicalModels}, where we presented a type of
graphical model, i.e., Bayesian networks.

Once again, we would like to remind readers that the purpose of
the present chapter was not to present research results about \pradv
s. The data was mainly used for the purpose of illustration.  As
previously mentioned, more thorough studies on \pradv s were presented
in earlier publications of ours
\citep{Atayanetal2019,GastAtayan2019,Gastetal2019,Atayan2021b,Atayan2021a},
where some of the visualizations presented in this chapter were also
used.


%%% Local Variables:
%%% mode: latex
%%% TeX-master: "main"
%%% End:
